\section*{Bab \ref{ch:g10:common-ancestry} --- Pola Tubuh dan Nenek Moyang Bersama}

\begin{solution}{exo:g10:common-ancestry:1}
Simetri setangkup dengan ujung depan dan ujung belakang; rangka dalam
dengan lajur tulang belakang; tali saraf punggung di atasnya; saluran
pencernaan di sisi perut dengan jantung di bawahnya dekat ke depan
(dan tungkai atau sirip yang berpasangan).
\end{solution}

\begin{solution}{exo:g10:common-ancestry:2}
Organ dengan struktur dan kedudukan yang sama dalam \omterm{def:g10:common-ancestry:bodyplan}{pola tubuhnya}, apa
pun fungsinya. Tungkai depan tetrapoda: humerus, radius dan ulna,
tulang pergelangan, jari --- dalam sebuah lengan, sayap, sirip renang.
\end{solution}

\begin{solution}{exo:g10:common-ancestry:3}
Analog: fungsinya sama (terbang), strukturnya sama sekali berbeda ---
tulang, otot dan kulit pada burungnya, selaput rangka luar pada
kupu-kupunya. Keduanya tidak menyatakan apa pun tentang kekerabatan.
\end{solution}

\begin{solution}{exo:g10:common-ancestry:4}
Kadalnya: kotak terkecil yang memuat katak dan kadal adalah
``tetrapoda''; troutnya hanya berada dalam kotak ``\omterm{def:g10:common-ancestry:bodyplan}{vertebrata}'' yang lebih besar.
\end{solution}

\begin{solution}{exo:g10:common-ancestry:5}
Fosil dengan gabungan ciri yang antara; tahap bersama yang dilewati
embrio \omterm{def:g10:common-ancestry:bodyplan}{vertebrata}; kesesuaian perbandingan urutan molekuler dengan
pengelompokan anatominya.
\end{solution}

\begin{solution}{exo:g10:common-ancestry:6}
Mamalia (jejak rambut, susu; tungkai tetrapoda; bernapas dengan udara
dan, pada mamalia, perkembangan amniotik). Tulang panggulnya adalah
sisa tungkai belakang: paus berasal dari mamalia darat berkaki empat.
\end{solution}

\begin{solution}{exo:g10:common-ancestry:7}
Buayanya masuk ke ``amniota'' bersama kadal dan merpati (bukan ke
``mamalia''); salamandernya masuk ke ``tetrapoda'' di luar
``amniota'', bersama kataknya.
\end{solution}

\begin{solution}{exo:g10:common-ancestry:8}
Kelompok kekerabatan harus memuat seluruh keturunan moyang bersamanya.
Moyang yang sama bagi trout dan ikan paru juga menjadi moyang katak dan
setiap tetrapoda; kelompok yang memuat trout dan ikan paru tetapi
mengecualikan katak meninggalkan sebagian keturunannya. ``Ikan''
menamai sebuah cara hidup, bukan sebuah garis keturunan.
\end{solution}

\begin{solution}{exo:g10:common-ancestry:9}
Simpanse (0), mencit (9), ayam (13), katak (18), tuna (21), ragi (45):
mamalia, amniota, tetrapoda, \omterm{def:g10:common-ancestry:bodyplan}{vertebrata}, eukariot --- kotaknya
berurutan.
\end{solution}

\begin{solution}{exo:g10:common-ancestry:10}
Struktur yang tetap ada pada ikan dewasa muncul sekejap pada embrio
kita dan lalu diolah ulang: kantong insangnya menjadi bagian telinga,
tenggorokan dan kelenjar, ekornya terserap kembali. Embrio yang
melewati tahap moyang adalah apa yang diramalkan pewarisan dari moyang
bersama dan bukan apa yang diramalkan rancangan bagi sebuah tujuan.
\end{solution}

\begin{solution}{exo:g10:common-ancestry:11}
Jika burung berasal dari moyang menyerupai reptil, pasti ada moyang
yang punya ciri reptil dan ciri burung yang pertama sekaligus;
menemukan satu yang bertarikh di antara kedua kelompok itu adalah
ramalan yang terbenarkan. Tanpa nenek moyang bersama, makhluk yang
mencampur sifat dua macam yang terpisah tidak punya alasan untuk ada.
\end{solution}

\begin{solution}{exo:g10:common-ancestry:12}
Analog: fungsinya sama dan rancangan umumnya mirip, tetapi asal
embrionya berbeda dan strukturnya terbalik. \omterm{def:g10:common-ancestry:homology}{Homologi} akan bermakna
bahwa moyang bersama moluska dan \omterm{def:g10:common-ancestry:bodyplan}{vertebrata} sudah punya mata kamera;
analogi bermakna bahwa moyang itu paling banyak punya sebidang kecil
peka cahaya, dan kedua garis keturunannya membangun kameranya sendiri-sendiri.
\end{solution}

\begin{solution}{exo:g10:common-ancestry:13}
Manusia dan mencit berbagi moyang bersama yang lebih belakangan
daripada yang dibagi masing-masing dengan ayam; kedua garis
keturunannya telah terpisah dari garis keturunan ayam selama waktu yang
sama, sehingga keduanya menumpuk kira-kira sebanyak perubahan
terhadapnya. Jarak manusia--ayam dan mencit--ayam mengukur selang yang sama.
\end{solution}

\begin{solution}{exo:g10:common-ancestry:14}
Ular punya telur amniotik, rangka dan tengkorak kadal, dan sebagian
(ular sanca, ular piton) menyimpan sisa tungkai belakang dan sebuah
panggul; fosil ular berkaki kecil memang ada. ``Tetrapoda'' berarti
berasal dari moyang berempat tungkai, entah tungkainya tetap ada atau tidak.
\end{solution}

\begin{solution}{exo:g10:common-ancestry:15}
Cara hidup menjelaskan bentuk sebuah sirip renang, bukan mengapa sirip
renang paus memuat tulang sebuah lengan yang tersusun seperti tulang
kelelawar --- dayung yang rata tidak memerlukan pergelangan. Cara hidup
tidak menjelaskan kantong insang pada embrio manusia, yang tidak
berguna baginya. Dan cara hidup tidak menjelaskan mengapa \omterm{def:g10:chemistry-of-life:families}{protein}
pernapasan, yang tak terlihat oleh lingkungannya, berbeda antarspesies
persis sebanyak yang diramalkan rangkanya. Hanya pewarisan dari moyang
bersama yang menerangkan ketiganya.
\end{solution}

\begin{solution}{pb:g10:common-ancestry:1}
\textbf{1.} Lajur tulang belakangnya: \omterm{def:g10:common-ancestry:bodyplan}{vertebrata}.

\textbf{2.} Empat tungkai menetapkan katak, kadal, merpati, mencit:
tetrapoda. Telur amniotik menetapkan kadal, merpati, mencit: amniota.

\textbf{3.} Tidak: ciri yang hadir pada satu \omterm{def:g10:biodiversity-scales:species}{spesies} tidak
mengelompokkan apa pun. Ciri itu mengenali \omterm{def:g10:biodiversity-scales:species}{spesiesnya} (dan akan
menetapkan sebuah kelompok jika lebih banyak kerabat --- burung lain,
mamalia lain --- ditambahkan).

\textbf{4.} \omterm{def:g10:common-ancestry:bodyplan}{Vertebrata} $\supset$ tetrapoda $\supset$ amniota, dengan
troutnya di luar kotak kedua dan kataknya di luar kotak ketiga;
kadal, merpati dan mencitnya di dalam kotak ketiga.

\textbf{5.} Tungkai kataknya: tulang sama, urutan sama. Ia sepadan
dengan ciri ``empat tungkai'', yang tidak dimiliki siripnya.

\textbf{6.} Kadal dan merpati (12). Troutnya (19--21 dari semuanya).

\textbf{7.} Merpati (13), kadal (14), katak (18), trout (21).

\textbf{8.} Kadal (16), merpati (17), mencit (18), trout (19):
amniotanya berjarak kira-kira sama, troutnya di luar itu --- taat asas
dengan kotaknya (kataknya lebih dekat ke setiap amniota daripada ke troutnya).

\textbf{9.} Bahwa kadal dan merpati berpisah satu sama lain setelah
garis keturunannya berpisah dari garis keturunan mencit: di dalam
amniota, kadal dan merpati lebih dekat kerabatnya. Tabelnya tidak dapat
memberitahukannya: tidak ada ciri di dalamnya yang dibagi kadal dan merpati saja.

\textbf{10.} Garis keturunan troutnya berpisah dari garis keturunan
bersama keempat lainnya sebelum keempatnya saling berpisah:
masing-masing telah terpisah dari troutnya selama waktu yang sama,
sehingga masing-masing menumpuk kira-kira sebanyak beda.

\textbf{11.} Trout di luar semua yang lain; keempat tetrapodanya
bersama; ketiga amniotanya bersama di dalamnya.

\textbf{12.} Anatomi dan sebuah \omterm{def:g10:chemistry-of-life:families}{protein} pernapasan tidak punya alasan
untuk sepakat kecuali keduanya merekam sejarah yang sama. Dua saksi
yang saling bebas dan menceritakan kisah yang sama membuat keserupaan
yang kebetulan menjadi tidak masuk akal.

\textbf{13.} Ke dalam ``mamalia'', di samping mencitnya (rambut, susu,
dan jarak molekuler yang terkecil). Sayapnya adalah tungkai depan
tetrapoda yang diubah, analog dengan sayap merpati, bukan \omterm{def:g10:common-ancestry:homology}{homolog}
dengan terbang berbulu; molekulnya (13 dari merpati, 8 dari mencit) membenarkannya.

\textbf{14.} Sitokrom c berubah terlalu lambat untuk memisahkan kerabat
yang dekat, dan ia bukan \omterm{def:g10:universal-dna:gene}{gen} yang bertanggung jawab atas beda yang
kasatmata. Ia menerangkan kekerabatan yang dalam, bukan \omterm{def:g10:universal-dna:gene}{gen} yang
membuat sebuah \omterm{def:g10:biodiversity-scales:species}{spesies} menjadi dirinya.

\textbf{15.} Eukariot: \omterm{def:g10:cells-common-unit:cell}{sel} dengan \omterm{def:g10:cells-common-unit:organelle}{inti sel} dan \omterm{def:g10:cells-common-unit:organelle}{mitokondria}, yang
berbagi \omterm{def:g10:chemistry-of-life:families}{protein} pernapasan ini.

\textbf{16.} Ke-13 bedanya ditumpuk dua garis keturunan sepanjang
320 juta tahun: $13/2 \approx 6.5$ tiap garis keturunan, yaitu sekitar
2 tiap 100 juta tahun sepanjang satu garis keturunan.

\textbf{17.} Katak--mencit 18, katak--merpati 17: rata-ratanya 17.5,
yaitu 8.75 tiap garis keturunan; pada 2 tiap 100 juta tahun, sekitar
430 juta tahun.

\textbf{18.} Rata-rata troutnya sekitar 20: 10 tiap garis keturunan,
sekitar 500 juta tahun.

\textbf{19.} Perkiraannya (430, 500) lebih tua daripada tarikh fosilnya
(370, lebih dari 420) tetapi berorde benar dan berurutan benar; lajunya
tidak persis tetap dan hitungannya dibulatkan, sehingga kesesuaian
dalam 20\% sudah sebaik yang dibolehkan caranya.

\textbf{20.} Keduanya menegakkan kekerabatan bersarang yang sama ---
trout, lalu katak, lalu ketiga amniotanya --- sedangkan molekulnya
menambahkan sebuah jam: taksiran berapa lama lalu tiap cabangnya berpisah.
\end{solution}
